% journal_methods.tex — Sections 2-4 of the journal manuscript (paper.tex),
% kept in a separate file while the results sections are written.

\section{Surrogate-based propagation and the surrogate contribution}
\label{sec:framework}
\subsection{Propagation of distributions through a surrogate}
Let $Y=f(\mathbf X)$ be the measurement model, $g_{\mathbf X}$ the joint
probability density of the inputs (here with independent components) and
$\{(\mathbf x_i,y_i)\}_{i=1}^N$, $y_i=f(\mathbf x_i)$, the $N$ available
model runs. A surrogate $\hat f$ is fitted to the runs and the Monte Carlo
method of GUM~S1~\cite{jcgm101} is applied to $\hat f$: $M$ draws
$\mathbf x^{(r)}\sim g_{\mathbf X}$ give $\hat y^{(r)}=\hat f(\mathbf x^{(r)})$,
from which the estimate $y$ (the mean), the standard uncertainty $u(y)$ (the
standard deviation), the probabilistically symmetric 95\% coverage interval
$[y_{0.025},y_{0.975}]$, the 95th percentile $y_{0.95}$ and the probability
$P(Y>T)$ of exceeding a threshold $T$ are computed. We call these the
\emph{reported quantities}.

Surrogates are fitted on a \emph{latent} scale $z=\varphi(y)$, with
$\varphi=\log$ for the log-variable variants and the identity otherwise, and
predictions are mapped back by $\varphi^{-1}$. For a Gaussian latent
predictive distribution, $\exp$ of the latent mean is the predictive median
of $Y$; no bias correction is applied.

\subsection{The surrogate error as an uncertainty contribution}
\label{sec:contrib}
The surrogate error $e(\mathbf x)=\varphi(f(\mathbf x))-\varphi(\hat f(\mathbf x))$
is unknown. Following Kennedy and O'Hagan~\cite{kennedy2001}, who treat model
inadequacy as an uncertainty source, we propagate
\begin{equation}
Y=\varphi^{-1}\!\left(\varphi(\hat f(\mathbf X))+E\right),
\label{eq:contrib}
\end{equation}
with $E$ a random variable describing the surrogate error, drawn
independently for every Monte Carlo trial. Three models for $E$ are
compared, all available from the $N$ runs alone:
\begin{itemize}\setlength\itemsep{1pt}
\item \emph{GP}: $E\mid\mathbf x\sim\mathrm N(0,s^2(\mathbf x))$, with
$s(\mathbf x)$ the predictive SD of a Gaussian process surrogate;
\item \emph{CV-normal}: $E\sim\mathrm N(0,s_{\mathrm{LOO}}^2)$, with
$s_{\mathrm{LOO}}^2=N^{-1}\sum_i r_i^2$ and
$r_i=\varphi(y_i)-\varphi(\hat f_{-i}(\mathbf x_i))$ the leave-one-out (LOO)
residuals; $\hat f_{-i}$ is the surrogate re-fitted, hyperparameters
included, without run $i$;
\item \emph{resampling}: $E$ drawn with equal probabilities from
$\{r_1,\dots,r_N\}$.
\end{itemize}
The same ingredients give pointwise 95\% prediction intervals for
$f(\mathbf x)$: the GP interval $\hat z\pm1.96\,s(\mathbf x)$, the CV-normal
interval $\hat z\pm1.96\,s_{\mathrm{LOO}}$, and the jackknife+
interval~\cite{barber2021}
\begin{equation}
\left[\,q^-_{\alpha}\{\hat z_{-i}(\mathbf x)-|r_i|\},\;
q^+_{\alpha}\{\hat z_{-i}(\mathbf x)+|r_i|\}\,\right],
\end{equation}
where $q^-_\alpha$ is the $\lfloor\alpha(N+1)\rfloor$-th smallest and
$q^+_\alpha$ the $\lceil(1-\alpha)(N+1)\rceil$-th smallest of the $N$ values,
$\alpha=0.05$. For exchangeable runs the jackknife+ interval covers $f$ with
probability at least $1-2\alpha$ whatever the surrogate~\cite{barber2021};
it is a member of the conformal family~\cite{angelopoulos2023}.

\emph{Why resampling the LOO residuals can be expected to work.} For an
input $\mathbf X\sim g_{\mathbf X}$ drawn independently of the runs,
$\varphi(f(\mathbf X))=\varphi(\hat f(\mathbf X))+e(\mathbf X)$ exactly, so
equation~\eqref{eq:contrib} gives the true distribution of the output if $E$
behaves as $e(\mathbf X)$. If the run inputs and $\mathbf X$ are
exchangeable, as for independent draws from $g_{\mathbf X}$, each $r_i$ is
distributed as the error of a surrogate fitted to $N-1$ runs at a new input,
so the empirical distribution of the $r_i$ estimates the distribution of
$e(\mathbf X)$, on the conservative side because $N-1$ rather than $N$ runs
are used~\cite{arlot2010}. Three approximations remain. Resampling draws $E$
independently of $\mathbf X$: it reproduces the variance of the error but
drops its covariance with $\varphi(\hat f(\mathbf X))$, which is zero in
expectation when $\hat f$ is the posterior mean of a correctly specified
model but not in general. A Latin hypercube design is close to, not exactly,
exchangeable with independent draws. And a feature of $f$ that no run
reaches leaves no trace in any $r_i$ (Section~\ref{sec:res-robust}). At a
fixed input $\mathbf x$, the same construction with $\hat f_{-i}(\mathbf x)$
in place of $\hat f(\mathbf x)$ is the predictive distribution behind the
jackknife+ interval, one of the conformal predictive distributions, which
are valid in finite samples under exchangeability~\cite{barber2021,vovk2019}.
For the distribution of $Y$, which mixes over the inputs, no such guarantee
is available, and the coverage study of Section~\ref{sec:res-coverage} is
its empirical check.

For a reported quantity $\theta$ itself (the mean, SD, $y_{0.95}$ or
$P(Y>T)$ of the output over a fixed set of input draws), two 95\% intervals
are compared: the 2.5th and 97.5th percentiles of $\theta$ over joint sample
paths of the GP posterior, and the jackknife interval
$\hat\theta\pm1.96\,\mathrm{SE}_{\mathrm J}$ with
$\mathrm{SE}_{\mathrm J}^2=\frac{N-1}{N}\sum_i(\hat\theta_{-i}-\bar\theta)^2$,
where $\hat\theta_{-i}$ is computed from $\hat f_{-i}$.

\emph{Hyperparameter uncertainty.} The GP intervals above plug in the
maximum-likelihood (ML) hyperparameters, which may make them too narrow
with few runs. Three variants test whether the GP's shortfall is a plug-in
effect: (i) a Student-t process, with the overall variance scale and a
constant mean integrated out under the reference prior
$p(\beta,\sigma^2)\propto\sigma^{-2}$ and the correlation hyperparameters at
their ML values (the standard Bayesian emulator~\cite{ohagan2006}, $N-1$
degrees of freedom); (ii) a fully Bayesian GP, all hyperparameters
integrated out by Markov chain Monte Carlo (affine-invariant ensemble
sampler~\cite{foremanmackey2013}, uniform prior on each log-hyperparameter
within the ML search range, 800 steps of which 300 burn-in, 200 posterior
draws; $N=20$ only, for cost); and (iii) the GP SD recalibrated by the LOO
standardised residuals $z_i=r_i/s_{-i}(\mathbf x_i)$: pointwise
$\hat z\pm q\,s(\mathbf x)$ with $q$ the $\lceil0.95(N+1)\rceil$-th
smallest $|z_i|$, $E=s(\mathbf x)\,z_J$ in equation~\eqref{eq:contrib}, and
sample paths with the covariance scaled by the mean of $z_i^2$. The last is
a hybrid that keeps the GP's localisation of the error and takes its scale
from the runs. Each variant is assessed in the same three ways.

\section{Test models}
\label{sec:models}
Table~\ref{tab:models} lists the twelve models. Eight have independent uniform
inputs, three have the non-uniform input distributions of a GUM-type
evaluation, and one is a CFD simulation. The physics model is the benchmark
of the earlier analysis; its quadratic term carries practically all the
output variance on $[-200,200]^3$ and lies in the class of a degree-2
polynomial. The plume model is the ground-level concentration of a
continuous ammonia release from a Gaussian plume with ground reflection,
\begin{equation}
C=\frac{10^3\,Q}{2\uppi\,u\,\sigma_y\sigma_z}
\exp\!\left(-\frac{y^2}{2\sigma_y^2}\right)
2\exp\!\left(-\frac{H^2}{2\sigma_z^2}\right)\ \mathrm{mg\,m^{-3}},
\end{equation}
with Briggs rural coefficients for neutral stability,
$\sigma_y=0.08x(1+10^{-4}x)^{-1/2}$ and $\sigma_z=0.06x(1+1.5\times10^{-3}x)^{-1/2}$.
The borehole~\cite{harper1983,morris1993,vlse},
Ishigami~\cite{vlse}, OTL circuit and piston~\cite{benari2007} and wing
weight~\cite{forrester2008} functions are standard test functions for
computer experiments. To test whether the results hold when the inputs
outnumber what 20 runs can resolve, the OTL circuit is also run with 14
additional inputs that do not affect the output (uniform on $[100,200]$,
as in the screening study of Moon~\cite{moon2010}), so that $d=20$ and
$N=20$ is one run per input. The flood
model~\cite{iooss2015} gives the maximal annual underflow of a river with
truncated Gumbel, truncated normal, triangular and uniform inputs. The
borehole model is also run with the input model of Harper and
Gupta~\cite{harper1983} (normal borehole radius, lognormal radius of
influence) and the plume model with GUM-type inputs: lognormal emission rate
(median 500\,g\,s$^{-1}$, geometric SD 2), a Weibull wind speed (shape 2,
scale 4\,m\,s$^{-1}$, truncated to $[0.5,15]$\,m\,s$^{-1}$) and a normal
crosswind offset (SD 15\,m). All formulas and input models were checked
against the independent implementation in UQTestFuns~\cite{wicaksono2023}
(relative differences below $10^{-10}$ for outputs and $10^{-6}$ for
marginal quantiles).

A further analytical model, reported separately, tests the limits of the
procedure: a smooth
function of four uniform inputs on $[0,1]^4$ plus a narrow Gaussian bump,
$f=x_1+x_2^2+0.5\sin(2\uppi x_3)+0.5x_4+3.6\exp(-\lVert\mathbf x-0.7\rVert^2/(2\cdot0.13^2))$.
The bump carries about 10\% of the output variance in a region (bump above
10\% of its height) of 3\% of the input volume, which a Latin hypercube of
20 runs misses with probability 0.53 and one of 40 runs with probability
0.26 (Section~\ref{sec:res-robust}).
Two variants test assumptions of the main study. In the first, the inputs
are dependent: the borehole model with GUM-type inputs, with a Gaussian
copula of correlation 0.5 between the two transmissivities and between the
two heads. Designs are drawn in the space of independent uniform variables
and mapped to the inputs by the inverse Rosenblatt transformation of the
copula, so the surrogates and the propagation are otherwise unchanged. In
the second, the output is a vector: the plume model at three receptors on
the plume axis, at $x$, $x+200$\,m and $x+400$\,m, gives three strongly
correlated output quantities, propagated as in GUM~S2~\cite{jcgm102}. Each
component has its own log-scale GP, and the surrogate error is added on the
log scale in one of three ways: as the LOO residual vector of one run, drawn
with equal probabilities (\emph{joint} resampling, which keeps the
correlation between the components), as residuals drawn independently for
each component, or as independent normal errors with the LOO RMS error of
each component. The coverage region is the hyperellipsoidal region of
GUM~S2, from the mean and covariance matrix of the propagated sample with
the coverage factor set so that 95\% of the sample lies inside, computed on
the original and on the log scale. Its achieved coverage is evaluated under
$2\times10^5$ draws of the true model, and its volume is compared with that
of the same region computed from the true model.

Exceedance is assessed at a fixed threshold where the model has one (the
IDLH level for ammonia, 210\,mg\,m$^{-3}$, about 300\,ppm, for the plume
and CFD models), and otherwise at the reference 95th percentile, for which
$P(Y>T)=0.05$ by construction. Flooding has a probability below $10^{-5}$
for the flood input model, so its threshold is the reference 5th percentile
of the underflow, $P(Y<T)=0.05$.

\begin{table*}[tb]
\centering
\caption{Test models. $d$: number of inputs; U: independent uniform inputs.
The models below the line test the limits and assumptions of the procedure
and are reported separately.}
\label{tab:models}
\small
\begin{tabular}{llp{53mm}p{54mm}}
\toprule
Model & $d$ & Inputs & Character\\
\midrule
Physics & 3 & U$[-200,200]^3$, noise SD 0.1 & quadratic bowl plus small oscillatory terms\\
Plume & 5 & U: $Q$, $u$, $H$, $x$, $y$ & multiplicative, positive, heavy right tail\\
Borehole & 8 & U, ranges of~\cite{vlse} & smooth, multiplicative, positive\\
Ishigami & 3 & U$[-\uppi,\uppi]^3$ & strongly non-linear, interactions\\
OTL circuit & 6 & U~\cite{benari2007} & smooth, nearly additive, positive\\
Piston & 7 & U~\cite{benari2007} & non-linear, positive\\
Wing weight & 10 & U~\cite{forrester2008} & smooth, multiplicative, positive\\
OTL circuit, 20 inputs & 20 & U, 14 inactive~\cite{moon2010} & as OTL circuit; sparse\\
Borehole (GUM) & 8 & normal, lognormal, uniform~\cite{harper1983} & as borehole; input tails\\
Plume (GUM) & 5 & lognormal, truncated Weibull, uniform, normal & as plume; heavier input and output tails\\
Flood & 8 & truncated Gumbel and normal, triangular, uniform~\cite{iooss2015} & smooth, lower-tail threshold\\
CFD dispersion & 6 & U: $U_{\mathrm{ref}}$, $z_0$, $H_{\mathrm b}$, $Sc_t$, $x_r$, $Q$ & RANS solution: wake recirculation, numerical error\\
\midrule
Bump (failure mode) & 4 & U$[0,1]^4$ & smooth part plus a narrow bump that the runs can miss\\
Borehole, dependent & 8 & as Borehole (GUM), Gaussian copula & dependent inputs\\
Plume, 3 receptors & 5 & as Plume & vector output (GUM~S2)\\
\bottomrule
\end{tabular}
\end{table*}

\subsection{CFD dispersion case}
\label{sec:cfd}
The CFD case replaces the plume formula by a Reynolds-averaged
Navier--Stokes (RANS) simulation of a continuous ground-level release
upwind of a building, solved with OpenFOAM v1912~\cite{weller1998}. The
domain is two-dimensional (a vertical plane along the wind; the building
and the release are infinitely wide), from 80\,m upstream to 250\,m
downstream of the windward face and 150\,m high. The building is 10\,m deep
and $H_{\mathrm b}$ high; the release is a unit line source 20\,m upstream
of the building, 0.5\,m above the ground. The inflow is the neutral
atmospheric surface layer of Richards and Hoxey~\cite{richards1993},
$U(z)=(u_*/\kappa)\ln((z+z_0)/z_0)$, $k=u_*^2/\sqrt{C_\mu}$,
$\varepsilon=u_*^3/(\kappa(z+z_0))$, scaled to the wind speed
$U_{\mathrm{ref}}$ at 10\,m, with the standard $k$--$\varepsilon$ model
($\sigma_\varepsilon=1.11$ for consistency with the inflow), a rough-wall
function on the ground~\cite{blocken2007}, smooth building walls and a slip
top boundary. The concentration $s$ of the passive scalar solves
$\nabla\cdot(\mathbf U s)-\nabla\cdot((\nu/Sc+\nu_t/Sc_t)\nabla s)=S$, with
$Sc=0.7$ and the turbulent Schmidt number $Sc_t$ as an input, since its
appropriate value is uncertain and varies between flows~\cite{tominaga2007}.
Flow and scalar are solved together by the SIMPLE algorithm on 7776
hexahedral cells, refined towards the building, until all scaled residuals
fall below $10^{-5}$ (pressure) and $10^{-6}$ (other fields), or 4000
iterations, typically in
1000--1500 iterations (about 20\,s on one core). The output is the
concentration at 1.5\,m height a distance $x_r$ downstream of the leeward
face, $C=10^3\,Q\,s(x_r)$\,mg\,m$^{-3}$ for a release rate $Q$ per metre of
crosswind width; since the scalar is passive, $x_r$ and $Q$ need no
separate run. The six inputs are uniform: $U_{\mathrm{ref}}\in[1,6]$\,m\,s$^{-1}$,
$z_0\in[0.01,0.3]$\,m, $H_{\mathrm b}\in[6,15]$\,m, $Sc_t\in[0.5,1.3]$,
$x_r\in[5,120]$\,m and $Q\in[1,20]$\,g\,s$^{-1}$\,m$^{-1}$.

The case contains what the plume formula lacks: flow separation at the
building edges, a recirculating wake whose length and strength change with
$H_{\mathrm b}$ and $z_0$, receptors inside and beyond the wake, and a
discretisation and iteration error that varies with the inputs. It is a
benchmark, not a validated prediction of a particular site: the model is
\emph{defined} as this discretised RANS solution, which is what a surrogate
of an expensive simulation has to reproduce. A reference Monte Carlo of
$2\times10^5$ CFD runs is not affordable. The reference distribution is
built instead from 800 runs at independent uniform draws of the four flow
inputs, each combined with 50 independent draws of $(x_r,Q)$: 40\,000
reference points, of which 4000 (five per run) form the test set and 1600
(two per run) the input set for derived quantities. The grid sensitivity of
the case is reported in Section~\ref{sec:res-cfd}.

\section{Surrogates, designs and protocol}
\label{sec:protocol}
\subsection{Surrogates}
\paragraph{Gaussian processes.} All GPs are fitted on standardised inputs
and outputs by maximising the log marginal likelihood (L-BFGS-B, five
random restarts) with scikit-learn. The specifications, in order of
increasing completeness, are: (i) the \emph{original} kernel of the earlier
analysis, a squared exponential with one length scale and an amplitude
\emph{fixed} at 1 on the standardised output; (ii) with the amplitude
estimated; (iii) with one length scale per input (automatic relevance
determination, ARD) but the amplitude fixed; (iv) with both; (v) a
Mat\'ern-5/2 ARD kernel with amplitude and a nugget; (vi) (v) plus an
amplitude-scaled quadratic polynomial kernel, i.e.\ a GP with a quadratic
trend whose coefficients are integrated out. Length scales are bounded to
$[10^{-2},10^3]$, the amplitude to $[10^{-3},10^3]$ and the nugget variance
to $[10^{-8},10^{-1}]$. Variants (v) and (vi) are also fitted on the log
scale for positive outputs, with logarithmic transformation of inputs whose
support is positive.

\paragraph{Sparse polynomial chaos.} A sparse PCE~\cite{blatman2011} with an
orthonormal Legendre basis on inputs mapped to $[-1,1]$ (affinely for
uniform inputs, through their distribution functions otherwise), candidate
sets of total degree $p=1,\dots,5$ with hyperbolic truncation ($q=1$ for
$p\le2$, $q=0.75$ above), least-angle regression~\cite{efron2004} to order
the basis terms, ordinary least squares on every active set along the path,
and selection of the active set and degree by the corrected LOO
error~\cite{blatman2011}. It is fitted on the raw and on the log scale.

\paragraph{Committees.} Nine members: two GPs (squared exponential and
Mat\'ern-3/2, both ARD with amplitude), PCEs~\cite{xiu2002,sudret2008} of degrees 2
and 3 (Bayesian ridge regression), random forests of 50 and 100 trees, a
support vector regressor and two neural networks. Members are combined
uniformly or with weights
$w_k\propto[\mathrm{MSE}_{\mathrm{CV},k}(1+\pi_k)]^{-1}$, with
$\mathrm{MSE}_{\mathrm{CV},k}$ the three-fold cross-validation error of
member $k$ and $\pi_k\ge0$ a penalty for predictions below a known
physical bound~\cite{goel2007,viana2009}.

\paragraph{Stacking.} As an alternative to choosing one candidate surrogate
(Section~\ref{sec:res-select}), the candidates are combined by
stacking~\cite{breiman1996}: non-negative weights that minimise the squared
error of the weighted sum of their LOO predictions (non-negative least
squares) are applied to the candidates fitted to all runs.

\subsection{Designs}
Passive designs are $N$-point Latin hypercubes drawn on $[0,1]^d$ and mapped
to the inputs by the inverse distribution functions. Adaptive designs start
from $N_0$ Latin hypercube runs ($N_0=5$ for $N\le20$, $N/4$ above) and add
one run at a time, refitting after every run: \emph{GP-SD} adds the
candidate with the largest posterior SD~\cite{gramacy2009}; the adaptive
committees add the candidate where the active members disagree most, by the
variance of their predictions or by the spectral measure
UEigV~\cite{lin2024} of the earlier analysis; and \emph{AK-MCS} adds the
candidate minimising $|\hat z(\mathbf x)-\varphi(T)|/s(\mathbf x)$, the
learning function of Echard et al.~\cite{echard2011}, which concentrates
runs near the threshold $T$ of the exceedance probability. \emph{IMSE}
adds the candidate that most reduces the integrated posterior variance of
the GP~\cite{sacks1989}, averaged over 1000 draws from the input
distribution, a criterion aimed at the whole distribution of the output
rather than at the largest local uncertainty. Candidates are random draws
from the input distribution (300 for GP-SD and IMSE, 2000 for AK-MCS).
To test whether the results depend on the type of passive design, the
Latin hypercube is replaced, in a separate comparison at $N=20$, by the
maximin Latin hypercube~\cite{morris1995} (the one with the largest
smallest distance between two runs among 200 random ones) and by the first
$N$ points of a randomly scrambled Sobol' sequence~\cite{sobol1967}.

\subsection{Protocol and statistics}
Every strategy is run with the 30 seeds 1000--1029 and budgets
$N=10,20,40,80$ (the full set of strategies at $N=20$). Seeds are shared and
the Latin hypercube is seeded from the global random state, so that for a
given seed all passive strategies use the identical design and all adaptive
strategies the identical start. Accuracy over the input domain is the root
mean square error on 5000 test points divided by the reference output SD
(normalised RMSE, nRMSE; 1 corresponds to a constant predictor). The
surrogates propagate $3\times10^4$ input draws; the reference uses
$2\times10^5$ draws of the model (Section~\ref{sec:cfd} for the CFD case).

For each strategy and model the median over seeds is reported with a 95\%
percentile-bootstrap confidence interval (10\,000 resamples of the seeds).
Two strategies are compared on the same seeds by the median of the per-seed
ratio of their errors, with a paired bootstrap 95\% confidence interval,
and by a two-sided Wilcoxon signed-rank test on the log ratios; the
$p$-values of all comparisons of the main study are adjusted together by
Holm's method, and those of each further comparison (design types, IMSE,
stacking) within that comparison. Coverage of prediction and coverage intervals is the share of test
points or of the reference distribution inside the interval, averaged over
seeds; coverage of intervals for reported quantities is the share of seeds
whose interval contains the true value, with a Wilson 95\% interval. The
interval score~\cite{gneiting2007} combines width and coverage. Because a
practitioner has a single design, coverage is also reported per design, as
the share of the 30 designs whose achieved coverage is below 90\%.

A comparison without a significant difference does not show that two
strategies are equivalent. Such comparisons are classified as equivalent
within $\pm20\%$ (or $\pm10\%$) when the 95\% bootstrap interval of the median
ratio lies inside $[0.8,1.25]$ (or $[0.9,1/0.9]$), and as inconclusive
otherwise. These margins, the thresholds for coverage (at least 93\%
pointwise; at least 26 of 30 seeds for a reported quantity, the smallest
count not rejected against 95\% by a one-sided exact binomial test at the
5\% level) and the decisions that the results would trigger were fixed in a
research plan committed to the repository~\repocite\ before the runs.
Finally, the sparse PCE and GP implementations were checked against an
independent one, OpenTURNS~1.27~\cite{baudin2017}, on the 30 designs of six
models with uniform inputs.
